\subsection{Bounded kernels, finite-algebra characters, and an actual word relation}
Allowing bounded complexes does not by itself repair the negative
word trace. The decisive comparison is whether the actual algebra
relations hold on full cohomology or only after taking Euler classes.

\begin{finding}[\fid{F60}: bounded-kernel relation gates]
\label{f:boundedkernelcharactergate}
Let \(A\) be a finite-dimensional unital algebra over a
characteristic-zero field \(F\). For every representation
\[
\rho:A\longrightarrow M_m(F[z,z^{-1}]),
\]
the character \(\operatorname{Tr}\rho(a)\) of each fixed
\(a\in A\) is independent of \(z\).
No semisimplicity assumption is needed.

Consequently there is no bounded internally graded five-state
matrix-kernel lift of the actual \(N^\#\) module in \fid{F57}
whose full cohomology-character matrices represent the actual
generic \(\check H_4\), with its scalar grading parameter unchanged.
In particular all actual relations cannot lift as ordinary
isomorphisms with scalar grading objects in cohomological degree zero.

Even the one rank-one module relation of the specified word,
in the precise isomorphism form below, is already impossible.
These conclusions do not apply to relations asserted only as
triangles or Euler-class equalities.
\end{finding}

\begin{proof}
\par\smallskip\noindent\textbf{Finite algebra, constant character.}\quad
For \(a\in A\), finite dimensionality gives a nonzero polynomial
\(p_a(T)\in F[T]\) with \(p_a(a)=0\), by linear dependence of
\(1,a,a^2,\ldots\).
Choose a finite constant extension \(E/F\) splitting this polynomial.
The matrix \(\rho(a)\) over \(E(z)\) is annihilated by \(p_a\),
so each eigenvalue over an algebraic closure of \(E(z)\) is one of
the fixed roots of \(p_a\) in \(E\).
Its trace is an integer multiplicity sum of those constants and
therefore lies in \(E\). Jordan blocks do not affect the trace.
But the trace also lies in \(F[z,z^{-1}]\).
Linear independence of the powers of the indeterminate \(z\)
over \(E\) makes every nonconstant Laurent coefficient zero.
The constant coefficient lies in \(F\). Thus
\begin{equation}\label{eq:finitealgebracharconstant}
\operatorname{Tr}\rho(a)(z)=\operatorname{Tr}\rho(a)(-1).
\end{equation}
Repeated roots of \(p_a\) cause no change to this proof.

There is an optional stronger conclusion when \(A\) becomes
split semisimple over a finite extension \(E\).
For each central unit \(e_r\) of a matrix factor \(M_{d_r}(E)\),
\(\rho(e_r)\) is idempotent. Its trace is its integer rank
\(n_r\), both generically and at \(z=-1\).
The matrix units identify its \(d_r\) diagonal image spaces,
so the corresponding simple multiplicity is \(n_r/d_r\).
These multiplicities agree at the two parameters.
Thus the modules become isomorphic after the splitting extension.
Without semisimplicity only the character conclusion is asserted.

\par\smallskip\noindent\textbf{The full cohomology character.}\quad
For a bounded complex \(C\) of finite internally graded
\(\Bbbk\)-vector spaces, with internal-degree-zero differential, set
\[
h_C(v,z)=\sum_{e,j}\dim_{\Bbbk}H^j(C)_e\,v^ez^j
\in\mathbb N[v,v^{-1},z,z^{-1}].
\]
The Euler character is \(h_C(v,-1)\); total cohomology has
character \(h_C(v,1)\).
All sums are finite.
Choosing graded complements to boundaries in cycles and to
cycles in terms splits a bounded complex into cohomology and
contractible pairs. Tensoring a contractible pair remains
contractible by its contracting homotopy. This proves the finite
K\"unneth formula in the form used here.

For matrix kernels use
\[
(K\star L)_{ab}=\bigoplus_c K_{ac}\otimes_{\Bbbk}L_{cb}.
\]
K\"unneth and finite direct sums give the exact matrix identity
\begin{equation}\label{eq:kernelPoincareproduct}
h_{K\star L}(v,z)=h_K(v,z)h_L(v,z).
\end{equation}
Derived isomorphisms, and in particular homotopy isomorphisms,
preserve this character.
Scalar finite graded objects in cohomological degree zero
have characters independent of \(z\).

\par\smallskip\noindent\textbf{The relation hypothesis.}\quad
Let \(g_1,\ldots,g_r\) generate \(A\) over \(F=\Q(v)\).
The explicit requirement is that
\[
g_i\longmapsto h_{K_i}(v,z)
\]
factors through an algebra representation
\begin{equation}\label{eq:fullcohomologyaction}
A\longrightarrow M_m(F[z,z^{-1}]),
\end{equation}
and that at \(z=-1\) it is the prescribed actual Euler module.
This follows if enough actual relations lift as ordinary
isomorphisms of sums of word kernels with cohomologically
degree-zero scalar grading objects.
Indeed clear rational-function denominators, expand Laurent
coefficients, and move negative scalar terms to the opposite
side. Positive multiplicities and internal shifts give the
two sums. Their isomorphism and
\eqref{eq:kernelPoincareproduct} imply the full relation before
specializing \(z\).

A distinguished triangle generally gives only its Euler
equality, not equality of full cohomology characters.
A coefficient object with nonzero cohomological degrees
can likewise change the scalar before \(z=-1\).
Neither is being assumed to satisfy
\eqref{eq:fullcohomologyaction}.

For every named word \(w\),
\eqref{eq:finitealgebracharconstant} now gives
\[
\operatorname{Tr}h_{K_w}(v,z)
=\operatorname{Tr}_{M}(g_w).
\]
The left side is a nonnegative Laurent polynomial in \(v,z\).
Evaluation at \(z=1\) therefore forces
\begin{equation}\label{eq:boundedkernelpositiveword}
\operatorname{Tr}_{M}(g_w)\in\mathbb N[v,v^{-1}].
\end{equation}
In fact its nonzero-cohomological-degree coefficients all vanish,
because the character is independent of \(z\).

IV, \S11.8, p.~67, gives the actual rank-four algebra's
dimension \(114\) and its module list.
Proposition~11.8, p.~61, additionally gives semisimplicity and
a finite splitting extension for the optional same-module
interpretation \cite{IV}.
The trace argument requires only finite dimensionality.
For \(M=N^\#\), the actual word in \fid{F57} has negative
coefficients at \(v^{18}\) and \(v^{-18}\), contradicting
\eqref{eq:boundedkernelpositiveword}.

The full actual relation hypothesis matters.
IV, \S11.8, p.~67, explicitly says that the quadratic and
adjacent nonstandard braid relations do not generate the
rank-four relation ideal. Those smaller-rank checks alone
do not justify \eqref{eq:fullcohomologyaction}.

\par\smallskip\noindent\textbf{One explicit module relation.}\quad
Let \(W\) be the actual \fid{F57} word matrix on \(N^\#\),
and put
\[
\pi(v)=p((v+v^{-1})^2),
\]
with \(p\) as in \eqref{eq:rankfournegativepolynomial}.
The original \(Q\)-word contains the rank-one factor
\(Q_1Q_3=ue_2^t\). Hence its rank is at most one,
and its nonzero trace makes the rank exactly one.
Transposition preserves this.
Writing a rank-one matrix as \(ab^t\) gives
\[
W^2=a(b^ta)b^t=\operatorname{Tr}(W)W=\pi(v)W.
\]
This is a relation on the specified module, not an assertion
that the same equation holds in the whole algebra.

Write \(\pi=\pi_+-\pi_-\), where
\(\pi_-=v^{18}+v^{-18}\) and \(\pi_+\) comprises the
remaining positive terms of
\eqref{eq:rankfournegativeLaurent}.
Let \(E_\pm\) be finite internally graded vector spaces
in cohomological degree zero with these characters.
A proposed lift \(K_W\) of the word cannot satisfy
\begin{equation}\label{eq:rankonewordisomorphismgate}
K_W\star K_W\ \oplus\ (E_-\otimes K_W)
\simeq E_+\otimes K_W
\end{equation}
while retaining its nonzero Euler trace \(\pi\).

Indeed put \(C(v,z)=h_{K_W}(v,z)\).
The isomorphism would give \(C^2=\pi(v)C\).
At \((v,z)=(3,1)\), the finite matrix \(T=C(3,1)\)
has nonnegative real entries, whereas
\[
\pi(3)=-\frac{99834597608550400}{387420489}<0.
\]
Thus \(T^2=\pi(3)T\) forces every entry of \(T\) to
vanish: each left entry is nonnegative and each right
entry is nonpositive, strictly negative if the corresponding
entry of \(T\) is positive.
A finite Laurent polynomial with nonnegative coefficients
vanishes at \((3,1)\) only when every coefficient is zero.
All cohomology entries of \(K_W\) would therefore vanish,
contradicting its Euler trace \(\pi\ne0\).
\end{proof}

The last test needs only its displayed module-annihilator
isomorphism, not the full relation ideal or semisimplicity.
Neither test uses arithmetic purity.
They leave triangle-only relations, cohomologically nontrivial
coefficient lifts, and geometric convolution without this
matrix-cohomology character functor outside their scope.
Unbounded constructions and larger modules without an effective
specified-summand comparison are also not excluded.
An algebra action does not automatically lift every annihilator
relation of its Euler module as an ordinary isomorphism.
The ordinary Soergel route of \fid{F59} is not being identified
with a finite-state vector-space kernel category.

% Separate paragraph for Part III.
\paragraph{\fid{F60}'s cohomology character in Grothendieck's language.}
The internal grading is a split-torus action; \(z\) records
cohomological degree and is not an arithmetic Frobenius parameter.
Effective cohomology gives \(h_C(v,z)\), while \(z=-1\)
gives the Euler class.
The framework is SGA~3, I, 4.7.2--3 and cote~90, scan~17;
coefficient representation and \(\lambda/\sigma\) language
uses SGA~6, V, \S\S2--3 and cote~133, scan~18
\cite{SGA3,SGA6,Archive}.
No \(\lambda\)-ring structure on the noncommutative algebra
is asserted.
The bounded coefficient complexes lie in SGA~6, I,
Definition~4.7, pp.~121--122, and Example~4.8, p.~122.
The finite K\"unneth computation used here was proved explicitly.

The representation family is an affine finite free module on
\(\operatorname{Spec}F[z,z^{-1}]\) carrying the constant algebra.
The affine dictionary is EGA~I, Theorem~1.3.7, p.~86,
Corollary~1.3.8, p.~87, and Theorem~1.4.1, pp.~90--92
\cite{EGA}.
Every algebra element satisfies a constant polynomial, so its
trace is a sum of fixed roots and cannot vary with \(z\).
In the additional split semisimple case, central idempotents
define constant-rank summand bundles and determine the simple
multiplicities.
The one-word obstruction uses the actual rank-one module equation
and the effective cohomology cone at a positive character value.
These are proved comparison obstructions, not assertions from an
unpublished archive leaf or identifications of a canonical basis.
